% Additive manuscript fragment.  See INTEGRATION.md before importing.
% Requires the existing theorem environment and amsmath support.
% Insert the companion bibliography-items.tex inside the existing bibliography.
\section{Uniform Gaussian proofs and signed-kernel evaluation}
\label{gcc:section}
The following proof upgrades establish the weight-five Gaussian relation and
the five weight-six double reductions displayed above.  The remaining $S_4$
and triple-polylogarithm candidates retain their stated status.

\subsection{A two-row certificate at weight five}
The same-point iterated-integral shuffle gives, with $p+q=w$,
\begin{equation}\label{gcc:shuffle}
 \operatorname{Li}_p(i)\operatorname{Li}_q(i)
 =\sum_{a=1}^{w-1}
 \left[\binom{a-1}{p-1}+\binom{a-1}{q-1}\right]
 \operatorname{Li}_{a,w-a}(i,1).
\end{equation}
Out-of-range binomial coefficients are zero. Taking imaginary parts at
$(p,q)=(1,4)$ and $(2,3)$ yields
\begin{align}
 g_{1,4}+g_{2,3}+g_{3,2}+2g_{4,1}
 &=-\tfrac12\beta(4)\log2-\tfrac{7\pi^5}{46080},\label{gcc:R14}\\
 g_{2,3}+3g_{3,2}+6g_{4,1}
 &=-\tfrac{3}{32}G\zeta(3)-\tfrac{\pi^5}{1536}.\label{gcc:R23}
\end{align}
Multiplying the first equation by $960$ and subtracting $224$ times the second
proves the previously numerical weight-five identity exactly.  The same rows
give
\begin{align}
 g_{2,3}&=-6g_{4,1}-3g_{3,2}-\frac{\pi^5}{1536}
                         -\frac{3G\zeta(3)}{32},\label{gcc:g23}\\
 g_{1,4}&=4g_{4,1}+2g_{3,2}+\frac{23\pi^5}{46080}
       +\frac{3G\zeta(3)}{32}-\frac12\beta(4)\log2.
\end{align}
Thus at most two coordinates remain among these four Gaussian values, modulo
the displayed depth-one products.  Their numerical independence is not asserted.

\subsection{A parity identity proving the even-weight reductions}
Put $F_{a,b}(z)=\operatorname{Li}_{a,b}(z,1)$ and $w=a+b$. For
$z=e^{i\theta}$, $0<\theta<2\pi$, define
\[
 \mathcal B_k(z)=\frac{(2\pi i)^k}{k!}B_k\!\left(\frac\theta{2\pi}\right).
\]
\begin{theorem}[Inner-argument-one parity specialization]
\label{gcc:parity}
For positive integers $a,b$,
\begin{align}
 F_{a,b}(z)-(-1)^wF_{a,b}(\bar z)
 ={}&\sum_{\mu=b+1}^w(-1)^{b+\mu}\binom{\mu-1}{b-1}
       \zeta(\mu)\mathcal B_{w-\mu}(z)-\operatorname{Li}_w(z)\notag\\
 &+(-1)^a\sum_{\mu=a}^w\binom{\mu-1}{a-1}
       \operatorname{Li}_\mu(\bar z)\mathcal B_{w-\mu}(z).
\end{align}
\end{theorem}
\begin{proof}
In Panzer's increasing-index convention \cite{gcc:Panzer}, substitute
$(n_1,n_2;z_1,z_2)=(b,a;y,z)$ in equation (3.2) and let $y\to1$ through a
branch-consistent sector. The $\mu=b$ term and the separate subtraction combine
to $\operatorname{Li}_b(y)[\mathcal B_a(zy)-\mathcal B_a(z)]$. This tends to zero
also for $b=1$, since the bracket is $O(|y-1|)$ and
$\operatorname{Li}_1(y)=-\log(1-y)$. The remaining inner polylogarithms have index
at least two and tend to their zeta values.  This proves the formula without
introducing $\zeta(1)$.
\end{proof}
For $z=i$ and even $w$, divide the imaginary part by two and use
$\operatorname{Li}_n(i)=-2^{-n}\eta(n)+i\beta(n)$. This proves all five of the
weight-six Gaussian double reductions in their existing normalization, and
supplies every even weight by finite rational arithmetic. The continuation
\cite{gcc:continuation} includes 64 exact formulas through weight 16 and
independent rational interval comparisons for all of them.

\subsection{Analytic structure and certification}
The continuation also constructs an integrable signed density
$\kappa_{a,b}$ with moments $H_{n-1}^{(b)}/n^a$, zero mass, and one strict
negative-to-positive sign change. It represents
\[
 F_{a,b}(z)=\int_0^1\frac{z}{1-zu}\kappa_{a,b}(u)\,du.
\]
For every integer $a\ge1$ and real $b>0$, this proves that
$\operatorname{Im}F_{a,b}(e^{i\theta})$ has exactly one zero in $0<\theta<\pi$,
that the zero is simple and below $\pi/2$, and that $g_{a,b}<0$.

For integer indices, write $f(n)=H_{2n}^{(b)}/(2n+1)^a$ and
$Df(n)=f(n)-f(n+1)$. Then
\[
 E_N=\sum_{k=0}^{N-1}\frac{D^kf(0)}{2^{k+1}},\qquad
 E_N-C_b2^{-N}\le g_{a,b}\le E_N,
\]
where $C_1=1$ and $C_b=b/(b-1)$ for $b\ge2$. Every endpoint is rational.
The article proves the signed-kernel theorem, a uniform four-term zero
asymptotic, and the sharp Euler error; the certificate implementation is
additive and does not rely on the historical numerical scripts.

Finally, the still-conjectural $S_4$ evaluation is equivalent, using
\eqref{gcc:g23}, to
\[
 S_4\stackrel{?}{=}\frac{58}{7}g_{4,1}+\frac{24}{7}g_{3,2}
                  +\frac{19\pi^5}{3584}-2\beta(4)\log2.
\]
The two sides differ by at most $10^{-118}$ in the delivered rational
certificate, but equality is not proved. A separating vector in the
continuation shows that a specified weight-five depth-two product relation
system alone does not prove this identity. Neither that obstruction nor
numerical agreement is an independence or nonexistence theorem.
